\begin{afterword}
\summaryhead

The author imagines drawing objects from a barrel: 11 blue eggs and 8 red cubes. Asked to guess the color of the next egg-shaped object, a careful reasoner says ``blue'' and knows it is a guess from a small sample. But on seeing a tiger, nobody reasons like this; the brain infers ``dangerous'' at once. Once the objects are named ``bleggs'' and ``rubes,'' the inference feels like no inference at all.

So, the post argues, you cannot define a word any way you like. That would be true only if words were logical classes from which nothing new is inferred, but the brain categorizes and infers regardless of our stated rules. The Aristotelians held that human mortality was established by definition, yet went on recognizing people as human from their faces and bodies. Their wrong picture of their own thinking did not stop their perception, but mistakes of this kind can spoil the deliberate thinking we use to check first impressions, for example by making us careless in choosing definitions.

\responsehead

The post's claim is that naming a category sets off inferences whether or not we approve of them, so that choosing a word is never a neutral act. The claim is plausible, and it has more support than the post gives it. The post cites no study; research on infants, such as Waxman and Markow's (1995), finds that new words direct attention to object categories from about the first birthday.

The historical frame is looser than the claim it serves. Aristotle is made the author of rules under which ``no one could conclude Socrates was `human' until after he died.'' That follows only if mortality is part of the definition of man. Aristotle's own examples define man as a two-footed or walking animal and count mortality as a property; the definition with ``mortal'' in it is Porphyry's, adopted by later logicians. Aristotle's account of how we come to know kinds is also closer to the post than the post allows: perception of a particular already presents ``man,'' and first premises come by induction. What does fit the post is the certainty Aristotle claimed for the knowledge that results. The individual example is also off: Socrates died fifteen years before Aristotle was born, so Aristotle never classified him by sight. The point survives, because the post moves at once to living people.

The conclusion states the purpose of the sequence: a wrong theory of one's own thinking ``can severely mess up'' deliberate thought. The post's only historical case shows the other half of its claim, a wrong theory that did no harm. Examples of harm are left to later posts, which is reasonable for an opening post but leaves this one to assert its main practical claim.

In short: a plausible and supported claim about words and inference, set in a history that fits later Aristotelians better than Aristotle, with its practical conclusion left for later posts to show.
\end{afterword}
